% AI-written, AI cross-reviewed; no human sign-off; not a claim that RH is proved
% Atlas unit T19. Local manuscript v3, item 5.29, heading on printed PDF p. 46.
% UpperTheory.LiteratureTheorems.resolution_char_zero_of_literature
% paper2/split_p012.tex:337–348. Verbatim excerpt; SHA-256 3b47d11a7accf3d44677ee217d781f8f4ddad5e05450bb0024436a9037ae4ca0.
% Original TeX uses the macros and cross-reference keys of the source manuscript.


% Start of atlas edition annotation. The following block is outside the original excerpt.
\paragraph{Status of named inputs in this reading.}
\noindent\nolinkurl{PositivityTheory.Resplug.ResolutionBlowup.resolutionCharZero} --- in-repo constant (\path{Development/WeilBound/CoreStatements/ResolutionCharZero/ResolutionCharZero.lean:6--7}); not compared.\par
% End of atlas edition annotation.

\begin{formalstatement}[Resolution over a characteristic-zero field]\label{sub:T19}

Given the open-literature record, for every characteristic-zero field $k$
and reduced separated finite-type $k$-scheme $X$, there is a sequence of
blowups $f:X'\to X$ with $X'$ smooth. Its restriction over the smooth
locus is an isomorphism. The reduced inverse image of the complement
is a simple normal crossings divisor. The formal wrapper uses a closed
resolution term independently of its record argument.

\end{formalstatement}\subnode{T19}

\noindent Kollár's Theorem 36(1)--(3) gives the three geometric conclusions. In Lean, \nolinkurl{IsBlowupSeq} starts with the identity and composes blowups characterized by the invertibility of the pulled-back ideal and its universal factorization. The predicate \nolinkurl{FinalTheorems.Snc} asks locally for an equation given by the product of distinct members of a regular system of parameters. These are the local equations of the reduced inverse image in Kollár's Definition 24. See \path{Geometry₂/ConclusionOfSmooth.lean:17--22} and \path{RegularLocalRings/FinalTheorems.lean:6--21}. \href{https://arxiv.org/pdf/math/0508332}{J. Kollár, Resolution of Singularities — Seattle Lecture, Theorem 36(1–3)}~\cite{r21-kollar}.
